\section{Computational Cost}
\label{sec:computational-cost}

\method{} does not add extra model forward or backward passes. For a fixed task, model, and epoch budget, the model-level FLOPs are therefore unchanged. The relevant additional cost is wall-clock overhead from controller updates, axis management, gradient-noise operations, logging, and other implementation details.

To estimate this overhead, each configuration was run three times with the same random seed on the same RTX PRO 6000 GPU. Runtime was measured at the epoch level, starting immediately before the training pass and ending after the corresponding validation evaluation. The first epoch was excluded from the timing average to reduce one-time warmup effects from CUDA initialization, dataloader startup, caching, and related setup costs. The reported values therefore reflect steady-state epoch time. They include both training and validation time, but exclude one-time setup costs such as dataset construction, tokenization, dataloader creation, and first-epoch warmup.

The repeated timing runs are not intended to estimate generalization variability. They are used only to estimate runtime overhead under fixed experimental conditions. Since validation does not involve \method{} controller updates or gradient-noise operations, the reported overhead should be interpreted as end-to-end steady-state epoch overhead rather than pure training-step overhead.

\begin{table}[t]
  \centering
  \caption{
    Wall-clock training overhead of scheduler and \method{} configurations, measured on a single NVIDIA RTX PRO 6000.
    Values are averaged over repeated timing runs with the same random seed.
    The \emph{Baseline} column reports seconds per epoch for the corresponding no-scheduler optimizer baseline;
    all remaining columns report percentage change in epoch wall-clock time relative to that baseline.
    \method{}-LR adapts only the learning-rate multiplier, \method{}-Noise adapts only gradient-noise magnitude, and \method{}-Dual adapts both axes.
  }
  \label{tab:compute-overhead}
  \sisetup{
    table-format = +2.1,
    table-space-text-post = {\,\%},
    detect-weight,
    round-mode = places,
    round-precision = 1
  }
  \begin{tabular}{
      l
      S[table-format=3.2]
      S[table-format=+2.1]
      S[table-format=+2.1]
      S[table-format=+2.1]
      S[table-format=+2.1]
    }
    \toprule
    & & \multicolumn{4}{c}{Overhead relative to baseline (\%)} \\
    \cmidrule(l){3-6}
    {Task / Optimizer}
      & {Baseline (s/ep)}
      & {LR sched.}
      & {\method{}-LR}
      & {\method{}-Noise}
      & {\method{}-Dual} \\
    \midrule
    CIFAR-100 / SGD & 5.17 & -0.0 & -0.5 & +3.4 & +4.0 \\
    CIFAR-100 / AdamW & 4.58 & +0.5 & +15.5 & +17.3 & +18.3 \\
    SST-2 / AdamW & 48.88 & +0.5 & +30.3 & +39.1 & +38.7 \\
    AG News / AdamW & 117.14 & +0.1 & +19.4 & +26.7 & +24.5 \\
    \bottomrule
  \end{tabular}
\end{table}

\Cref{tab:compute-overhead} shows that conventional learning-rate scheduling adds negligible overhead. In contrast, \method{} introduces non-negligible wall-clock cost, and the magnitude of this cost depends on the active control axis and training regime. For \cifar{} with \sgdm{}, LR-only \method{} adds almost no overhead, while noise-enabled variants add a few percent. For \cifar{} with \adamw{}, the overhead is larger, with LR-only control adding $15.5\%$ and dual-axis control adding $18.3\%$. The overhead is highest in the Transformer fine-tuning experiments, where LR-only control adds $30.3\%$ on \sst{} and $19.4\%$ on \agnews{}, while noise-enabled variants are generally more expensive.

Across the \adamw{} and Transformer settings, \method{}-LR is consistently cheaper than \method{}-Noise, and dual-axis control is close to the noise-enabled variants. This suggests that gradient-noise generation and application are major contributors to the additional runtime cost, while LR-only control is cheaper but still affected by implementation overhead.

Small negative overhead values, where present, should be interpreted as timing variability rather than meaningful speedups. These configurations do not reduce the underlying model computation, so apparent speedups mainly reflect measurement noise in short epoch timings.

The measured wall-clock overhead reflects the present implementation rather than an algorithmic lower bound. Part of the overhead comes from Python-level controller updates, axis bookkeeping, logging, and gradient-noise operations. A more optimized integration with the optimizer could reduce this cost, especially for LR-only control. Noise-enabled variants would still require additional work from gradient-noise generation and application.

These results qualify the accuracy improvements reported above. \method{} should not be described as computationally free, even though model FLOPs are unchanged for a fixed training budget. Its practical value depends on whether gains in peak validation accuracy, final validation accuracy, or stability justify the additional wall-clock time.

For transparency, the full per-configuration wall-clock and throughput measurements are reported in \Cref{tab:compute_overhead_detailed}. The detailed table is supplementary runtime evidence rather than a basis for averaging across heterogeneous \method{} variants.